\chapter{Instruments: Making Structure Visible to Bisimulation}
\label{ch:obsext}

The previous chapter generated a logic from a presentation. This one settles a
question about that logic which, left open, would quietly undermine everything the
next part builds on it --- and which a reader who knows the concurrency literature
will raise on sight.

\section{Two claims that look inconsistent}
\label{ox:sec:contradiction}

A generated logic is worth generating only if it is \emph{adequate}: logical
equivalence should coincide with the calculus's observational equivalence.
Chapter~\ref{ch:hypercube} was careful to say that adequacy is a property of the
logic one asks for rather than of the generator, and that a logic missing
conjunction or negation will not have it. What that section did not address is a
sharper objection, which does not turn on missing connectives at all.

The generated logic has a \emph{structural} layer: one connective per term
constructor, with
\[
  P \sat C(\varphi_1,\dots,\varphi_n)
  \iff
  P =_{\eqs} C(t_1,\dots,t_n) \text{ with } t_i \sat \varphi_i
  \text{ for each } i .
\]
When the constructor in question is the interaction cut, and the cut is subject to
associativity, commutativity and a unit, this connective is a \emph{separating
conjunction}: $P \sat \varphi \Par \psi$ holds exactly when $P$ admits
\emph{some} decomposition into two parts satisfying $\varphi$ and $\psi$
respectively. And that existential over decompositions is not a step. No modality
expresses it, because nothing happened.

So one wants to assert both of
\begin{align}
  \text{logical equivalence} \;&=\; \text{observational equivalence},
    \tag{A}\label{ox:eq:A}\\
  \text{the structural layer} \;&\subsetneq\; \text{observational equivalence}.
    \tag{B}\label{ox:eq:B}
\end{align}
Read carelessly these are inconsistent, and a careful reader will say so. The book
asserts both. \eqref{ox:eq:A} is what Chapter~\ref{ch:hypercube} wants and what the
Wigner argument of Chapter~\ref{ch:sci} leans on. \eqref{ox:eq:B} appears in
Chapter~\ref{ch:sci} too, where the strict refinement of context-labelled
bisimulation by the spatial layer is offered as the reason the ecology needs a
narrower class of morphisms than $\catGSLT$ supplies.

A referee put it to me exactly this way, and was right to. If structural
inspection is an admissible observation, then interaction was never the whole
observational ceiling; and if it is not, those predicates should not be among the
scientist's senses. One cannot have it both ways without saying more.

\section{The resolution}
\label{ox:sec:resolution}

There is no such thing as \emph{the} observational equivalence of a calculus
independent of what the observer is permitted to do.

\eqref{ox:eq:B} is a statement about bisimilarity in a theory $\GSLT$ where the
observer may build contexts out of $\GSLT$'s own constructors and watch what
happens. \eqref{ox:eq:A} is a statement about bisimilarity in a theory where the
observer additionally holds instruments for taking terms apart. These are two
members of a family. Exhibit the index and there is nothing to reconcile.

The construction that exhibits the index is the subject of this chapter. Its
slogan:

\begin{quote}
\itshape
Make every term former an interaction. Bisimulation observes interactions.
Therefore bisimulation observes every term former.
\end{quote}

The phenomenon is not new. Caires observed the apparent gain in discriminating
power in the spatial logic for the $\pi$-calculus~\cite{Caires,CairesCardelli},
and within a few years results in that community established that suitable
behavioural equivalences see what the spatial connectives see, so that spatial
logics need not be counted intensional~\cite{Lozes,HirschkoffLozes}. Those results are semantic: they exhibit
a calculus, show that its observational equivalence is characterised by a spatial
logic, and the extra separating power comes from features of that calculus. What
follows is instead \emph{syntactic}. It is a transformation of presentations: read
the signature, adjoin some atoms, add three rules per constructor, stop. It applies
to any interactive GSLT meeting the conditions of \S\ref{ox:sec:prelim} rather than
to one calculus, and --- this is the part that matters here --- because it is
mechanical in the presentation, it composes with the machinery of
Chapter~\ref{ch:hypercube}, which is also mechanical in the presentation.
Adequacy then becomes the statement that two functions of the same argument agree,
rather than a coincidence to be checked calculus by calculus.

\begin{remark}[Three things to keep apart]
\label{ox:rem:three}
An earlier version of this material braided three claims together tightly enough
that a weakness in one looked like a weakness in all three. They are:
\emph{operational reconstruction}, what the enlarged transition system can
distinguish (\S\ref{ox:sec:recon}); \emph{logical expressiveness}, what the
generated structural language can distinguish (\S\ref{ox:sec:sep}); and
\emph{observer capability}, which instruments a given observer actually holds
(\S\ref{ox:sec:dial}, \S\ref{ox:sec:security}). They are proved separately below
and should be cited separately.
\end{remark}

\begin{remark}[What is not being claimed]
\label{ox:rem:notclaim}
It is \emph{not} being claimed that structural inspection was interaction all
along. That would be slippery. The form of observation remains transition
observation; what changes is that the set of observable transitions has been
enlarged, and the enlargement is exactly what supplies the additional
discriminating power. The defensible statement is that structural observation
\emph{can be represented as} interaction, after conservatively extending the
observer's transition theory with destructuring actions. An architecture that
appeals to this construction should say its observer holds instruments, not that
its instruments were never instruments.
\end{remark}

\section{What is fixed, and what is a parameter}
\label{ox:sec:prelim}

Fix an interactive GSLT $\GSLT = (\Sigma, \eqs, \rules)$ in the sense of
Chapter~\ref{ch:igslt}: a signature, a set of equations, a set of rewrite rules,
and a distinguished binary constructor $\Cut$ --- the cut --- such that the
left-hand side of every rule is of the form $\Cut(l_1,l_2)$ up to $=_{\eqs}$.
Parallel composition plays this role for CCS, the $\pi$-calculus, the ambient
calculus and \rhoc; application plays it for the $\lambda$-calculus.

\subsection{Labels are minimal enabling contexts}

Chapter~\ref{sec:lts-hml} fixed the label discipline: a transition
$P \xrightarrow{K} Q$ is asserted when $K$ is a \emph{least} context enabling a
rewrite. Two refinements of that chapter's presentation are used here and are
stated in Remarks~\ref{lts:rem:universal} and~\ref{lts:rem:higherorder} there:
minimality is a universal property rather than a subterm condition, and labels
carrying process payloads are matched up to the bisimilarity being defined rather
than syntactically.

Minimality is not a technical convenience. It is what keeps a label from carrying
the whole term. A transition system labelled by \emph{arbitrary} enabling contexts
would be so informative that matching labels would nearly force syntactic
agreement, and the equivalence it induced would already sit at or near $=_{\eqs}$
--- collapsing the very distinction this chapter exists to draw. We return to this
in \S\ref{ox:sec:nonvac}, where it does real work.

\begin{condition}[Redex RPOs]
\label{ox:cond:rpo}
$\GSLT$ has all redex-relative pushouts, so that the derived transition system
exists.
\end{condition}

\begin{condition}[Congruence]
\label{ox:cond:cong}
$\GSLT$ comes equipped with a class $\mathcal{A}$ of contexts, containing those
built from the rewrite constructors, such that context-labelled bisimilarity is a
congruence with respect to $\mathcal{A}$, and observations are restricted to
$\mathcal{A}$.
\end{condition}

Both are inherited from~\cite{BIHOL}, where establishing them for the theories of
interest is work in progress; everything below is conditional on them. They are
conditions rather than theorems deliberately. We are interested only in
context-labelled transition graphs that make bisimulation a congruence. Barbed
bisimulation, which obtains congruence by quantifying over contexts and a
calculus-specific notion of observable, does not generalise across interactive
GSLTs and is not used. Deriving labels from contexts so that congruence holds by
construction is the whole point of the Leifer--Milner
apparatus~\cite{leifer2000deriving}.

Condition~\ref{ox:cond:cong} is not idle. For \rhoc, congruence holds under
$\mathrm{out}(n,-)$, $\mathrm{in}(n,-)$, $-\Par-$ and $\ast(@(-))$, and must
exclude contexts with a hole \emph{under} the quote, such as
$\mathrm{out}(@(-),z)$: communication turns on name equality, not on
bisimilarity, so $p \bisim q$ with $p \neq q$ gives $@p \neq @q$ and the two
contexts behave differently. \S\ref{ox:sec:security} returns to what this costs
the present construction.

\subsection{Admissible equational theories}

The results need $\eqs$ to admit a well-founded notion of decomposition. Better to
isolate the conditions than to discover them mid-proof.

\begin{definition}[Measure]
\label{ox:def:measure}
For a term $P$ let $|P|$ be the number of constructor occurrences and
$\mu(P) = \min\{\,|Q| : Q =_{\eqs} P\,\}$. A representative attaining the minimum
is \emph{lean}.
\end{definition}

\begin{definition}[Admissible $\eqs$]
\label{ox:def:adm}
$\eqs$ is \emph{admissible} when:
\begin{description}
\item[(E1) Well-founded.] Every $=_{\eqs}$-class of closed terms contains a lean
representative.
\item[(E2) Hereditary.] If $C(\vec t\,)$ is lean then each $t_i$ is lean, and for
$n \ge 1$ we have $\mu(t_i) < \mu(C(\vec t\,))$.
\item[(E3) Sorted.] $\eqs$ relates only terms of the same sort.
\end{description}
\end{definition}

\begin{remark}[Why (E2) is not free]
\label{ox:rem:unit}
It is tempting to say instead that $\eqs$ is size-respecting on the presentations
of interest. It is not. The unit law $P \Par \mathbf{0} = P$ is not size-respecting,
and it is one of the equations doing the most work below, since it is what makes
$P \Par \mathbf{0}$ one of the decompositions the separating conjunction quantifies
over. (E1)--(E2) are the repair. They say nothing about how many decompositions a
term has --- possibly infinitely many --- only that one may always descend along a
lean one. Associativity, commutativity and unit for the cut satisfy them.
\end{remark}

\section{The construction}
\label{ox:sec:construction}

The cut is binary; constructors have arbitrary arity and heterogeneous argument
sorts. To present the arguments of an $n$-ary constructor to the cut they must be
gathered into one operand.

\begin{remark}[Not lists, and not Turing completeness]
\label{ox:rem:noturing}
One might encode the arguments as a list and appeal to Turing completeness to
supply the list former. That would be wrong twice over. Computability guarantees
the ability to represent data; it does not supply an internal coding function
injective modulo an arbitrary $\eqs$, and injectivity modulo $=_{\eqs}$ is exactly
what the reconstruction below needs. And a general lambda theory is many-sorted,
so a homogeneous list cannot hold heterogeneous arguments without a universal code
sort, which is a further construction choice. The repair is simpler than the thing
repaired: \emph{freely adjoin} an argument former per signature profile. Free
adjunction gives correct sorting and injectivity by construction, introduces no
encoding choice, and keeps Turing completeness out of the chapter entirely.
\end{remark}

\begin{definition}[Observer extension]
\label{ox:def:ext}
Let $\Ob \subseteq \Sigma$ be an \emph{observation signature}. The \emph{observer
extension} $\Obs(\GSLT,\Ob) = (\Sigma^{\Ob}, \eqs, \rules^{\Ob})$ is given by:
\begin{itemize}
\item for each $C \in \Ob$ of profile
  $\tau_1 \times \cdots \times \tau_n \to \tau$, a fresh sort $A_C$ and a freely
  adjoined former $\args_C : \tau_1 \times \cdots \times \tau_n \to A_C$;
\item nullary \emph{probe atoms} $\Ask_C$, $\Bld_C$ and $\Get_{C,i}$ for
  $1 \le i \le n$;
\item the cut overloaded so that $\Cut(\Ask_C, t)$, $\Cut(\Get_{C,i}, a)$ and
  $\Cut(\Bld_C, a)$ are well formed for $t : \tau$ and $a : A_C$;
\item $\eqs$ unchanged;
\item $\rules^{\Ob} = \rules \cup \{\,\mathrm{ask}_C,\ \mathrm{get}_{C,i},\
  \mathrm{bld}_C : C \in \Ob\,\}$ where
\begin{align*}
  \mathrm{ask}_C &: \;\;
    \Cut(\Ask_C,\, C(x_1,\dots,x_n)) \;\rewrite\; \args_C(x_1,\dots,x_n),\\
  \mathrm{get}_{C,i} &: \;\;
    \Cut(\Get_{C,i},\, \args_C(x_1,\dots,x_n)) \;\rewrite\; x_i,\\
  \mathrm{bld}_C &: \;\;
    \Cut(\Bld_C,\, \args_C(x_1,\dots,x_n)) \;\rewrite\; C(x_1,\dots,x_n).
\end{align*}
\end{itemize}
Rules of $\rules$ are \emph{proper}; the rest are \emph{administrative}. Write
$\GSLT^{+}$ for $\Obs(\GSLT,\Sigma)$.
\end{definition}

Four comments on the shape of this definition, because each records a choice that
could have gone otherwise.

\paragraph{The cut is overloaded rather than a fresh eliminator introduced.}
This is the point of the whole construction. The cut is not an arbitrary
constructor; it is the one the transition relation is indexed on, the one that
defines what it means for two things to meet. A probe placed in interaction
position turns a structural question into a rendezvous, and a rendezvous is what a
derived transition system is built to observe.

\paragraph{Every administrative rule is interaction-headed by construction.}
Its left-hand side is literally a cut. Putting the constructor being opened on the
left of the arrow instead would force an appeal to a unit for the cut, and a
special case for cuts without one. Putting the probe on the left removes the case
distinction: the observer supplies the context, so no unit is needed anywhere.

\paragraph{The probe atoms are indexed.}
It is $\Ask_C$ rather than a single $\Ask$ that makes the label of an opening
transition record \emph{which} constructor was found, and $\Get_{C,i}$ rather than
a single $\Get$ that makes projection record \emph{which} argument was taken. That
is what carries the information the reconstruction needs.

\paragraph{The build rule is used in no proof below.}
The rule $\mathrm{bld}_C$ is present so that the extension supplies construction as
well as deconstruction, and so that the observer's instruments form a coherent kit rather
than a one-way mirror.

\begin{proposition}[The extension is still an interactive GSLT]
\label{ox:prop:still}
$\Obs(\GSLT,\Ob)$ is a GSLT, is interactive with the same cut, and reflects proper
transitions: for $\Sigma$-terms $P,P'$ and $\rho \in \rules$, $P \xrightarrow{K} P'$
in $\GSLT$ iff the same holds in $\Obs(\GSLT,\Ob)$.
\end{proposition}

\begin{proof}
Every rule of $\rules^{\Ob}$ has a left-hand side of the form $\Cut(-,-)$: those of
$\rules$ by hypothesis, the administrative ones by construction. Reflection of
proper transitions holds because $\rules^{\Ob} \supseteq \rules$ and the adjoined
formers are mentioned in no rule of $\rules$, so no $\rules$-redex is created or
destroyed by their presence, and the relative pushouts computed for
$\rules$-redexes are unchanged.
\end{proof}

\begin{remark}[Equations, or rewrite pairs]
\label{ox:rem:pairs}
One may prefer a presentation with no equations, orienting each equation of
$\eqs$ in both directions and adding both to the administrative rules. The two
presentations agree up to administrative steps. We keep $\eqs$ as equations
because doing so localises the source of nondeterminism: with $\eqs$ retained, the
branching of $\mathrm{ask}_C$ comes from $\eqs$-matching, which is exactly the
existential over decompositions that the separating conjunction quantifies. It is
worth stressing that these must be \emph{pairs}. Orienting the equations in one
direction only would let bisimilarity fall below $=_{\eqs}$ onto raw syntax, which
overshoots: the structural connectives respect $=_{\eqs}$, so a relation finer than
$=_{\eqs}$ is finer than the logic, and adequacy fails from the other side.
\end{remark}

\section{Reconstruction}
\label{ox:sec:recon}

\begin{lemma}[Exposure]
\label{ox:lem:exposure}
Let $P$ be a $\Sigma^{\Ob}$-term and $C \in \Ob$ be $n$-ary. Then
\[
  P \xrightarrow{\;\Cut(\Ask_C,\,-)\;} \args_C(t_1,\dots,t_n)
  \quad\text{iff}\quad
  P =_{\eqs} C(t_1,\dots,t_n),
\]
and $\Cut(\Ask_C,-)$ is a least context enabling $\mathrm{ask}_C$ at $P$.
Consequently the $\Ask_C$-labelled transitions of $P$ enumerate exactly the
top-level $\eqs$-decompositions of $P$ at $C$. Likewise $\args_C(\vec t\,)$ has the
single transition $\xrightarrow{\Cut(\Get_{C,i},-)} t_i$ with that label.
\end{lemma}

\begin{proof}
The redex $\Cut(\Ask_C, C(\vec x))$ is formed from $P$ by the context
$\Cut(\Ask_C,-)$ exactly when $P =_{\eqs} C(\vec t\,)$; no smaller context enables
it, since the rule's left-hand side requires the probe, and no larger one is
minimal. Matching is up to $=_{\eqs}$, which is where the enumeration comes from.
For $\mathrm{get}$, $\args_C$ is freely adjoined, so
$\args_C(\vec t\,) =_{\eqs} \args_C(\vec s\,)$ iff $t_i =_{\eqs} s_i$ for all $i$,
and the projection is deterministic.
\end{proof}

This is the lemma that does the work, and it is worth pausing on what it says in
the case that motivated the difficulty. Take $C$ to be the cut with $\eqs$
containing associativity, commutativity and unit. Then $P$'s $\Ask_{\Cut}$-labelled
transitions are in bijection with the ways of writing $P$ as $Q \Par R$ ---
including $P \Par \mathbf{0}$ and $\mathbf{0} \Par P$, and every rebracketing and
permutation of a parallel product. \textbf{The existential over splits in
$P \sat \varphi \Par \psi$ has become an existential over transitions.} Everything
else in this chapter is bookkeeping.

\begin{lemma}[Calibration]
\label{ox:lem:calib}
If $P =_{\eqs} Q$ then $P \bisim_{\Obs(\GSLT,\Ob)} Q$.
\end{lemma}

\begin{proof}
Transitions are defined up to $=_{\eqs}$, so $=_{\eqs}$-equal terms have identical
labelled transition sets, and the identity relation on $=_{\eqs}$-classes is a
bisimulation.
\end{proof}

Trivial as stated, and stated anyway, because it fixes the \emph{lower} bound of
the calibration: the enlargement must not push bisimilarity below the equational
theory, or it overshoots the logic. It is also the property against which the
alternative presentation of Remark~\ref{ox:rem:pairs} must be checked.

\begin{theorem}[Reconstruction]
\label{ox:thm:main}
Let $\GSLT$ satisfy Conditions~\ref{ox:cond:rpo} and~\ref{ox:cond:cong} with
$\eqs$ admissible, and let $\Ob = \Sigma$. Then for closed $\Sigma$-terms $P,Q$,
\[
  P \bisim_{\GSLT^{+}} Q \iff P =_{\eqs} Q .
\]
\end{theorem}

\begin{proof}
($\Leftarrow$) is Lemma~\ref{ox:lem:calib}.

($\Rightarrow$) By strong induction on $\mu(P)$, with induction hypothesis: for all
closed $P',Q'$ with $\mu(P') < \mu(P)$, if $P' \bisim_{\GSLT^{+}} Q'$ then
$P' =_{\eqs} Q'$. Note that the hypothesis constrains only the left argument, which
is what allows the right-hand witnesses produced below to be arbitrary
representatives rather than lean ones.

Let $P \bisim_{\GSLT^{+}} Q$ and choose a lean representative $C(t_1,\dots,t_n)$ of
$P$, which exists by (E1). By Lemma~\ref{ox:lem:exposure},
$P \xrightarrow{\Cut(\Ask_C,-)} \args_C(t_1,\dots,t_n)$. Since
$P \bisim_{\GSLT^{+}} Q$, the term $Q$ must match this step with a label related to
$\Cut(\Ask_C,-)$. That label has no process payload --- $\Ask_C$ is a nullary atom
and the hole is the only other position --- so label matching is identity on the
constructor skeleton, and $Q$'s matching label is $\Cut(\Ask_C,-)$ itself. Distinct
constructors carry distinct atoms, so no other opening rule could have matched. By
Lemma~\ref{ox:lem:exposure} again, $Q =_{\eqs} C(s_1,\dots,s_n)$ for some $\vec s$,
and $\args_C(\vec t\,) \bisim_{\GSLT^{+}} \args_C(\vec s\,)$.

Now apply $\Get_{C,i}$. By Lemma~\ref{ox:lem:exposure} each side has exactly one
transition with label $\Cut(\Get_{C,i},-)$, to $t_i$ and $s_i$ respectively, so
$t_i \bisim_{\GSLT^{+}} s_i$ for each $i$.

By (E2) each $t_i$ is lean and, for $n \ge 1$, $\mu(t_i) < \mu(P)$. The induction
hypothesis gives $t_i =_{\eqs} s_i$, whence
$P =_{\eqs} C(\vec t\,) =_{\eqs} C(\vec s\,) =_{\eqs} Q$. For $n = 0$ the
constructor is nullary and $P =_{\eqs} C =_{\eqs} Q$ directly.
\end{proof}

\begin{remark}[Domains]
\label{ox:rem:domains}
Theorem~\ref{ox:thm:main} is a statement about $\bisim_{\GSLT^{+}}$ restricted to
closed $\Sigma$-terms. The relation itself is defined on all of the extended
signature's terms, where it also identifies $\args_C(\vec t\,)$ with
$\args_C(\vec s\,)$ exactly when the components agree. The distinction matters in
Proposition~\ref{ox:prop:idem}.
\end{remark}

\begin{remark}[Infinite branching is not a problem]
\label{ox:rem:branching}
A term with replication has infinitely many $\eqs$-decompositions and therefore
infinitely many $\Ask_{\Cut}$-labelled transitions. Nothing in the proof requires
image-finiteness: the induction is on $\mu$, not on transition depth, and no
Hennessy--Milner style characterisation is invoked. \S\ref{ox:sec:sep} is what
makes this safe on the logical side as well.
\end{remark}

\begin{remark}[Binding, and the largest gap]
\label{ox:rem:binding}
For a binding constructor the arguments are abstractions rather than closed terms,
and Theorem~\ref{ox:thm:main} as proved does not cover them. It would need
$\alpha$-equivalence to be part of $\eqs$ and (E1)--(E2) to hold of the resulting
classes, a notion of transition beneath a binder, and a statement of what
$\Cut(\Ask_C,-)$ means when $C$ binds. The lambda-theoretic setting should supply
all three --- an abstraction is a term of the theory, so the induction has
somewhere to go --- but expecting is not proving. The theorem is therefore stated
for the first-order fragment. This is the largest gap between what this chapter
proves and what it would like to prove, and it is listed as such in the front
matter.
\end{remark}

\begin{proposition}[Idempotence]
\label{ox:prop:idem}
$\bisim_{\GSLT^{++}}$ and $\bisim_{\GSLT^{+}}$ agree on the terms of
$\Sigma^{\Sigma}$.
\end{proposition}

\begin{proof}
Since $\GSLT^{++}$ has more rules than $\GSLT^{+}$, monotonicity
(Proposition~\ref{ox:prop:mono}) gives
$\bisim_{\GSLT^{++}} \subseteq \bisim_{\GSLT^{+}}$. For the converse we exhibit a
bisimulation rather than invoke Theorem~\ref{ox:thm:main} on both sides, which
would be a domain mismatch of the kind Remark~\ref{ox:rem:domains} warns about.

Take $\mathcal{R} = \bisim_{\GSLT^{+}}$ and check the additional transitions. The
constructors adjoined by Definition~\ref{ox:def:ext} are of two kinds. The probe
atoms are nullary: their $\mathrm{ask}$ rules have targets with no components, so
the transition carries no information beyond the identity of the atom, which
$\bisim_{\GSLT^{+}}$ already fixes, since $\Cut(\Ask_C,-)$ appears in labels. The
argument formers are not nullary, but the $\Get_{C,i}$ rules of $\GSLT^{+}$ already
project every one of their components, so the second-round opening transition is
matched whenever the first-round projections are, and adds no separating power.
Hence $\bisim_{\GSLT^{+}}$ satisfies the $\GSLT^{++}$ matching conditions.
\end{proof}

This is not a side remark. It is the reason $\Obs$ is a construction rather than
the first stage of a colimit, and it turns on a design decision that could have
gone the other way: the probes are nullary and the argument formers are
exhaustively projected. Introduce a family of $n$-ary destructors with partial
projections instead and the ladder is real.

\section{Separation}
\label{ox:sec:sep}

The operational side is done. The logical side is a separate question. It would be
convenient to import it from the literature, but those are theorems about
particular spatial logics and particular calculi, and since the point here is to
generalise past the calculus-by-calculus treatment, borrowing the key step defeats
the purpose. We prove what is needed directly.

\begin{theorem}[Structural separation]
\label{ox:thm:sep}
Let $\eqs$ be admissible. Then for closed $\Sigma$-terms $P,Q$,
\[
  P =_{\mathcal{L}_{\mathrm{s}}(\GSLT)} Q \iff P =_{\eqs} Q ,
\]
where $\mathcal{L}_{\mathrm{s}}(\GSLT)$ is the structural fragment of the
generated logic.
\end{theorem}

\begin{proof}
($\Leftarrow$) Satisfaction of structural connectives is defined up to $=_{\eqs}$.

($\Rightarrow$) We construct a characteristic formula. For closed $P$, by (E1)
choose a lean representative $C(t_1,\dots,t_n)$ and set
\[
  \chi_P \;:=\; C(\chi_{t_1},\dots,\chi_{t_n}),
\]
well defined by induction on $\mu$, using (E2) for the descent, with
$\chi_P := C$ for nullary $C$.

We claim $Q \sat \chi_P$ iff $Q =_{\eqs} P$. If
$Q \sat C(\chi_{t_1},\dots,\chi_{t_n})$ then $Q =_{\eqs} C(s_1,\dots,s_n)$ with
$s_i \sat \chi_{t_i}$, so by induction $s_i =_{\eqs} t_i$, so
$Q =_{\eqs} C(\vec t\,) =_{\eqs} P$. Conversely $P \sat \chi_P$ by construction.
Hence if $P \neq_{\eqs} Q$ then $\chi_P$ separates them.
\end{proof}

\begin{remark}[What the direct proof buys]
\label{ox:rem:buys}
Three things. The argument is self-contained: no separation result is imported at
a generality its source does not support. The image-finiteness objection dissolves,
because the argument goes from the logic to $=_{\eqs}$ directly by characteristic
formulae and never passes through a Hennessy--Milner characterisation of
bisimulation, which is where image-finiteness or infinitary conjunction would have
been needed. And the hypotheses are visible: separation holds exactly for
admissible $\eqs$, so a presentation whose equational theory fails (E1)--(E2) is one
for which nothing is claimed. Note also that no boolean structure was used ---
$\chi_P$ is purely structural.
\end{remark}

\begin{corollary}[Logical correspondence]
\label{ox:cor:adequacy}
Under the hypotheses of Theorems~\ref{ox:thm:main} and~\ref{ox:thm:sep}, for closed
$\Sigma$-terms
\[
  =_{\mathcal{L}_{\mathrm{s}}(\GSLT)} \;=\; =_{\eqs} \;=\; \bisim_{\GSLT^{+}} .
\]
The structural fragment generated from $\GSLT$ is adequate for context-labelled
bisimulation in $\GSLT^{+}$.
\end{corollary}

That is the statement \eqref{ox:eq:A} should have been all along. It is adequacy
with an index.

\section{The gap is really there}
\label{ox:sec:nonvac}

Corollary~\ref{ox:cor:adequacy} would be worthless if bisimilarity in $\GSLT$ were
already $=_{\eqs}$, since the construction would then close a gap that was not
there. The same observation settles that and supplies \eqref{ox:eq:B} for
\emph{our} relation rather than the literature's.

\begin{proposition}[The gap is generically nonempty]
\label{ox:prop:nonvac}
Let $\GSLT$ contain two distinct nullary constructors $c,d$ of the same sort
occurring in no rule of $\rules$ and in no equation of $\eqs$. Then
$c \bisim_{\GSLT} d$ and $c \neq_{\eqs} d$, so
$\bisim_{\GSLT^{+}} \subsetneq \bisim_{\GSLT}$.
\end{proposition}

\begin{proof}
Neither $c$ nor $d$ occurs in any left-hand side, so neither contributes a redex,
and for any context $D$ the transitions of $D[c]$ and $D[d]$ are in
label-preserving bijection with those of $D[-]$'s own redexes. Hence
$\{(c,d)\} \cup {=_{\eqs}}$ extends to a bisimulation. They are $=_{\eqs}$-distinct
by hypothesis, so by Theorem~\ref{ox:thm:main} they are not
$\bisim_{\GSLT^{+}}$-related, witnessed by $\Cut(\Ask_c,-)$.
\end{proof}

The example is deliberately austere: it shows the gap is nonempty for essentially
any signature carrying inert data, without relying on a coincidence of the
equational theory.

\begin{example}[Replication, and a caveat]
\label{ox:ex:replication}
In a presentation with replication as a primitive constructor and $\eqs$
\emph{not} containing the unfolding law, $!P$ and $!P \Par P$ are behaviourally
indistinguishable and $=_{\eqs}$-distinct, so they witness the gap. If $\eqs$ does
contain the unfolding law they are $=_{\eqs}$-equal and witness nothing. The size
of the gap is a property of the presentation, not of the calculus informally
described --- a point worth carrying to any argument in this book that turns on
what a spatial connective can see.
\end{example}

\begin{remark}[Why minimality is load-bearing]
\label{ox:rem:minimality}
Proposition~\ref{ox:prop:nonvac} depends on labels being minimal. Were labels
arbitrary enabling contexts, a label would carry essentially the whole surrounding
term, matching would force near-syntactic agreement, and bisimilarity in $\GSLT$
would sit at or near $=_{\eqs}$ --- leaving no gap to close and no content to
Corollary~\ref{ox:cor:adequacy}. The label discipline of
Chapter~\ref{sec:lts-hml} is therefore not a technical preliminary but a
load-bearing part of the statement.
\end{remark}

\section{The dial}
\label{ox:sec:dial}

Nothing forces the construction to be run on the whole signature. The point of
Definition~\ref{ox:def:ext} taking $\Ob$ as a parameter is that the observer's
instruments are part of the data.

\begin{proposition}[Monotonicity]
\label{ox:prop:mono}
If $\Ob \subseteq \Ob'$ then
$\bisim_{\Obs(\GSLT,\Ob')} \subseteq \bisim_{\Obs(\GSLT,\Ob)}$, with
$\bisim_{\Obs(\GSLT,\varnothing)} = \bisim_{\GSLT}$ and, under the hypotheses of
Theorem~\ref{ox:thm:main}, $\bisim_{\Obs(\GSLT,\Sigma)} = {=_{\eqs}}$ on closed
$\Sigma$-terms.
\end{proposition}

\begin{proof}
The administrative rules in $\Obs(\GSLT,\Ob')$ but not $\Obs(\GSLT,\Ob)$ carry
fresh probe atoms, hence fresh labels, and add transitions that must be matched;
the largest relation satisfying more conditions is contained in the largest
relation satisfying fewer. The endpoints are Proposition~\ref{ox:prop:still} and
Theorem~\ref{ox:thm:main}.
\end{proof}

This is the sharp form of the resolution promised in \S\ref{ox:sec:resolution}.
There is no single observational bisimulation against which a spatial logic must be
measured. Observational equivalence is indexed by the destructuring capabilities
admitted to the observer, and \eqref{ox:eq:A} and \eqref{ox:eq:B} name the two ends
of the resulting family.

What does \emph{not} follow, tempting as it is, is that the fragment of the logic
retaining structural connectives only for $\Ob$ is adequate for
$\bisim_{\Obs(\GSLT,\Ob)}$.

\begin{conjecture}[Partial characterisation]
\label{ox:conj:partial}
For $\Ob \subseteq \Sigma$, $\bisim_{\Obs(\GSLT,\Ob)}$ coincides with logical
equivalence in the fragment of the generated logic whose structural connectives are
those of $\Ob$.
\end{conjecture}

The obstacle is real. A formula may inspect an opened constructor sitting
underneath an unopened one; equations may relate terms whose roots lie on opposite
sides of $\Ob$; and a binder may move material across an apparent structural
boundary. None of that is handled by relativising Theorem~\ref{ox:thm:main}.

\begin{proposition}[Downward closure]
\label{ox:prop:downward}
Conjecture~\ref{ox:conj:partial} holds when $\Ob$ is closed under the subterm
relation of lean representatives --- that is, when every constructor that can occur
beneath a constructor of $\Ob$ in a lean term is itself in $\Ob$.
\end{proposition}

\begin{proof}
Under that condition the descent in the proofs of Theorems~\ref{ox:thm:main}
and~\ref{ox:thm:sep} never leaves $\Ob$, so both arguments run verbatim with
$\Sigma$ replaced by $\Ob$ and $=_{\eqs}$ replaced by its restriction to the
$\Ob$-generated subalgebra.
\end{proof}

I suspect Conjecture~\ref{ox:conj:partial} is the most interesting statement in the
vicinity, and that its proof would say something about how partial structural
knowledge composes.

\begin{remark}[Three restrictions, kept apart]
\label{ox:rem:threerestrictions}
An agent observing its environment faces three different limits, and collapsing
them causes trouble downstream. There is the \emph{ideal observer} with the full
observation signature, whose reach is $=_{\eqs}$. There is the
\emph{capability-limited} observer, holding some $\Ob \subsetneq \Sigma$, whose
reach is the corresponding point of Proposition~\ref{ox:prop:mono}. And there is
the \emph{budget-limited} observer, who holds instruments it cannot afford to use,
whose reach is smaller again and is not characterised here at all. Adequacy is a
claim about the first. Chapter~\ref{ch:sci} distinguishes the first and the third
and will need the second as well.
\end{remark}

\section{Two cautions for anything built on this}
\label{ox:sec:cautions}

\subsection{Transitions and derivations}

Lemma~\ref{ox:lem:exposure} says the $\Ask_C$-labelled transitions of $P$ enumerate
its decompositions. The transition relation is extensional: $P \xrightarrow{K} P'$
holds when a witnessing $\eqs$-match exists. If two distinct $\eqs$-matches yield
the same label and the same target up to $=_{\eqs}$, they are one transition and
two derivations.

For bisimulation this is immaterial --- only existence is used. For any
quantitative lift it is not: weights and costs attach to derivations, and a
construction that collapses derivational multiplicity will misprice.
Chapter~\ref{ch:weight} is careful about this on its own account, taking a redex to
be a selection rather than a term up to congruence; anyone composing the two
constructions should carry the derivation structure explicitly rather than the
relation.

\subsection{Cost is bounded by depth, not equal to it}
\label{ox:sec:cost}

In a cost-accounted GSLT every rewrite carries a price, and administrative rules are
rewrites like any other. It is tempting to conclude that the price of evaluating a
structural formula of depth $d$ is a function of $d$, so that pricing schedules
stipulated elsewhere in this book become derived step counts. That is too quick.
Nesting depth bounds the length of the longest \emph{sequential} destructuring
path, but the cost of checking depends further on branching: a conjunction requires
two subchecks, a negation may require exhausting a search, and a separating
conjunction may enumerate many $\eqs$-decompositions before finding one that works,
or all of them before concluding that none does. Sharing and caching move the
figure again.

The defensible statement is that depth bounds sequential destructuring depth, while
actual checking cost depends additionally on branching and on the evaluation
strategy for the boolean layer. This is worth stating plainly because
Chapter~\ref{ch:sci} prices perception by formula depth and is explicit that the
schedule is a choice; the present chapter says why no other answer was available,
rather than leaving the choice looking like an oversight.

\section{The construction is not capability-safe}
\label{ox:sec:security}

This is stated plainly rather than hedged, because a reader who knows
object-capability discipline will see it in Definition~\ref{ox:def:ext} and will
want to know whether we noticed.

Unrestricted opening destroys encapsulation. Any observer may form
$\Cut(\Ask_C,-)$, so any term whatever may be taken apart, and $\Get_{C,i}$ then
yields the arguments as data. In a calculus where an opaque reference is the unit of
authority, this is fatal: open the reference and harvest whatever authority its
contents carry. Confinement fails, and with it every argument that depends on it.

The repair is known in outline. Mike Stay's suggestion is that constructors and
destructors should each take a capability token, so that the right to build and the
right to open are held rather than ambient. There are two readings --- a
\emph{permission}, licensing the rule, or a \emph{brand}, so that building stamps
the term and opening checks the stamp --- and the second is the more useful, since
it makes legibility a property of the term rather than a global property of the
world. The token must gate the whole kit rather than opening alone, since an
observer able to take an administrative step should be able to take its converse.

\begin{remark}[The factory is an import]
\label{ox:rem:import}
Remark~\ref{ox:rem:noturing} obtained the argument formers by free adjunction. That
cannot be done again here. Unforgeability is not a claim about what can be computed,
nor about what can be adjoined, but about what a context can \emph{write down}: in
a signature of freely adjoined constructors every term is writable by anybody, which
is precisely why free adjunction was safe for the argument formers and is useless
for a token. A genuine source of privacy must be imported --- restriction, in the
general case, giving a factory of the shape
$\mathsf{let}\ c = \mathsf{freshCap}()\ \mathsf{in}\ P$. Reflective calculi, in
which a name is a quoted process, may be able to discharge the import internally;
whether quotation \emph{originates} privacy or merely propagates it from a seed is
not settled here. Note that Chapter~\ref{ch:sci}'s freshness-by-quotation result
does not settle it either: not every interactive GSLT has a quote.
\end{remark}

\begin{remark}[Factories resist the construction]
\label{ox:rem:nodestructor}
Whatever the factory is, Definition~\ref{ox:def:ext} must not be applied to it.
Give it an opening rule and a token can be taken apart; if it can be taken apart it
can be rebuilt; if it can be rebuilt it can be forged; and if it can be forged
nothing is a capability. So the capability-safe version of the construction is not
``the observer extension, with a caveat'' but ``the observer extension on
$\Ob \subseteq \Sigma \setminus F$, where $F$ is the minting apparatus'' --- and the
exclusion is what makes the rest of it safe. It is tempting to say the factory is
the \emph{unique} constructor admitting no destructor. That is more than can be
supported: a system may carry several fresh-name generators, sealed abstract types,
or independent minting authorities, and in nominal presentations freshness is a
binder rather than a constructor at all. The property that matters is not
uniqueness but that the mechanism generating unforgeable authority is not itself
generically destructible.
\end{remark}

\subsection{Reflection, and an inherited restriction}

The congruence restriction of \S\ref{ox:sec:prelim} --- no hole beneath the quote
--- is inherited here and constrains $\Ob$. The opening rule for the quote,
$\Cut(\Ask_{@}, @(p)) \rewrite \args_{@}(p)$, does not itself place a hole beneath
$@$: the hole is beside the quote, not inside it. So one does not expect it to
disturb congruence. But it is a genuine addition to the context class, and whether
$\mathcal{A}$ extended by the administrative contexts still supports congruence is a
question for the machinery of~\cite{BIHOL} rather than for assertion here. A reader
wanting a safe default may take $\Ob$ to exclude the quote, at the cost of leaving
the name predicate outside the correspondence of
Corollary~\ref{ox:cor:adequacy}.

\section{What this settles, and where it is spent}
\label{ox:sec:settles}

\paragraph{Adequacy claims survive contact with intensionality results.}
A generated spatial--behavioural logic may be asserted adequate, provided the
assertion names the relation: bisimulation in the observer extension for a stated
$\Ob$, not in the base. Assertions that omit the index are not wrong so much as
ambiguous, and the ambiguity is what invites the objection of
\S\ref{ox:sec:contradiction}. Chapter~\ref{ch:hypercube}'s discussion of adequacy
should be read with this index supplied.

\paragraph{Spatial and behavioural addresses coincide.}
In $\Obs(\GSLT,\Ob)$ a decomposition site is a redex position, so constructions
identifying a structural locus with a modal label --- Chapter~\ref{ch:sci} has one,
where crossover loci are matched against modal labels --- are well typed rather than
analogical. This does \emph{not} supply an enabling theorem: a locus that is a redex
position for an administrative rule need not correspond to a \emph{proper} step
reachable in $\GSLT$, and claims about reachability still require their own
argument. The construction makes the identification meaningful; it does not make it
true.

\paragraph{Morphisms.}
Where the morphisms of a category of GSLTs preserve bisimulation
(Chapter~\ref{ch:catgslt}), an assignment depending on spatial structure is not
functorial, and the usual repair is to narrow the morphism class by hand to maps
preserving the generated logic. Under Theorem~\ref{ox:thm:main} the narrowing is not
a narrowing: logic-preserving maps of $\GSLT$ are exactly bisimulation-preserving
maps of $\GSLT^{+}$, and the repair becomes functoriality on the image of
$\Obs(-,\Sigma)$. Chapter~\ref{ch:sci} states this where it is needed.

\paragraph{And a warning.}
Example~\ref{ox:ex:replication} is the one to carry forward. How much the spatial
layer sees over and above bisimulation is a fact about the \emph{presentation}, not
about the calculus one has in mind. Two presentations of what a working programmer
would call the same language can sit at different points of
Proposition~\ref{ox:prop:mono}. Wherever this book argues from what a structural
predicate can detect, that argument is presentation-relative, and the presentation
is a modelling choice like any other.
